\chapter{Literature Survey}

This chapter includes summary of the findings/surveys that are carried as a ground work for the project. Findings/survey are to be from peer reviewed journals/conferences such as IEEE, IET, etc. .
\\ Sample  literature review in IEEE format \\
\textbf{1. Understand the Purpose:}\\
A literature review should:
Summarize and analyze existing research relevant to your project.\\
Identify trends, gaps, and challenges in the field.\\
Justify your research by showing how it builds on or differs from existing work.\\
\textbf{2. Structure of the Literature Review} \\
A well-structured IEEE literature review typically follows this format:\\
\textbf{A. Introduction}
Define the scope and purpose of the review.\\
Highlight the importance of the topic.\\
Outline the structure of the review.\\
\textbf{B. Review of Existing Work}\\
Organize this section thematically, chronologically, or methodologically: \\
Thematic Approach: Group related studies based on key topics or trends.\\
Chronological Approach: Arrange studies in the order of their publication to show project evolution.\\
Methodological Approach: Compare studies based on research methods used.
Each reviewed work should include:\\
Summary of the study (main objectives, methodology, key findings).\\
Relevance to your project.\\
Strengths and weaknesses of the study.\\
\textbf{C. Comparative Analysis and  Gap Identification}\\
Compare different approaches, findings, and methodologies.\\
Identify limitations in existing studies.\\
Highlight unresolved challenges and gaps in research.\\
\textbf{D. Summary and Justification for Your Project work}\\
Summarize key insights from the review.\\
Explain how your project builds on or addresses gaps in the literature.\\
Justify the need for your research.\\

\textbf{Sample review of Health Assistant system }
\\\textbf{A. Introduction}
With the increasing demand for digital health solutions, intelligent health assistant systems have emerged as an effective means to monitor patient health, provide recommendations, and facilitate doctor-patient interactions. These systems leverage artificial intelligence (AI), machine learning (ML), and the Internet of Things (IoT) to enhance healthcare services. This section reviews key research studies in the field, highlights advancements, and identifies gaps that justify the need for the proposed system.

\textbf{B. Review of Existing Work}
\textbf{AI-Based Health Assistants}
Recent studies have explored the use of AI-driven virtual health assistants for symptom checking and patient guidance. For example, in [1], a chatbot-based system was developed using Natural Language Processing (NLP) to provide preliminary diagnoses based on user-inputted symptoms. However, the study identified limitations in accuracy due to the dependency on predefined symptom databases.

\textbf{IoT-Based Health Monitoring Systems}
IoT-integrated health monitoring devices have gained popularity for real-time patient tracking. Work in [2] proposed a wearable IoT-based health assistant that continuously monitored heart rate, blood pressure, and oxygen levels. The system provided alerts to healthcare professionals in case of abnormalities. While effective, the system lacked predictive analytics for early disease detection.

\textbf{Machine Learning for Predictive Healthcare}
ML models have been extensively applied to predict health conditions based on historical patient data. The work in [3] utilized deep learning techniques to analyze patient records and predict chronic diseases such as diabetes and hypertension. Although the model achieved high accuracy, the lack of interpretability limited its real-world application.

\textbf{Mobile Health Applications}
Smartphone-based health assistants have revolutionized remote healthcare accessibility. In [4], a mobile application was developed to assist patients in tracking medication schedules, diet plans, and exercise routines. Despite its user-friendly design, the system faced adoption challenges among elderly users due to limited digital literacy.\\
\textbf{C. Comparative Analysis and Gap Identification}
Table I provides a comparative summary of key research contributions.
\begin{table}
    \centering
    \caption{Caption}
    \begin{tabular}{|c|c|c|c|} \hline 
        Study			 & Technology Used & Strengths  & Limitations\\ \hline
       [1]	  & AI Chatbot (NLP)		  &Automated symptom checking  &Limited accuracy  \\\hline
      [2]	  &IoT Sensors	   & Real-time health tracking	 & No predictive analytics \\\hline
       [3]	  & Deep Learning	 & Chronic disease prediction	 & Lack of interpretability \\ \hline
       [4] &	Mobile App	&User-friendly&	Low adoption among elderly \\ \hline
    \end{tabular}
    
    \label{tab:my_label}
\end{table}
While existing studies have significantly contributed to health assistant technologies, gaps remain in achieving high-accuracy, real-time, and predictive health assistance with better user adaptability.


\\\textbf{Summary and Justification for Proposed System}
\\This review highlights the need for an integrated AI-powered, IoT-enabled, and ML-driven health assistant that provides:
\\
Real-time health monitoring via IoT sensors.
\\Accurate and explainable AI-based diagnostics.
\\User-friendly mobile application for seamless interaction.
\\Predictive analytics for early disease detection.
\\The proposed system aims to address these challenges by combining IoT data, AI-driven diagnostics, and ML-based predictive analytics to offer an efficient and adaptive digital health assistant.

%This illustrates how to cite a reference. This is the first \cite{bib:paper} reference cited. To site a second \cite{bib:raju} reference, it's very easy. Refer bibliography.tex file to know , how to cite the reference.
